\section{Explicit two-witness colorings}\label{app:explicit}
This appendix gives a deterministic version of Lemma~\ref{lem:palette} on
selected domain sizes. The affine-plane construction follows the
direction-partition method of Jipsen, Maddux, and
Tuza~\citep{jipsen-maddux-tuza1995}. Their assignment of at least two
directions to every color already guarantees two witnesses at plane
order at least four; we include the verification because this is the
property used here.

\begin{proposition}[Explicit two-witness colorings]\label{prop:explicit-graphs}
For every palette size $N\ge2$, a graph with the two-witness property
of Lemma~\ref{lem:palette} can be constructed deterministically on fewer
than $16N^2$ vertices.
\end{proposition}
\begin{proof}
Let $s$ be the smallest power of two with $s\ge2N-1$, and let
$\mathbb F_s$ be the field of $s$ elements. Use vertex set $V\coloneqq\mathbb F_s^2$.
There are $s+1$ directions of affine lines. Partition these directions
into sets $\operatorname{Dir}_\mu$, $\mu\in\Lambda$, each of size at least
two. Color distinct points $u,v$ by the unique $\mu$ containing the
direction of their joining line. Since $N\ge2$, we have $s\ge4$ and
$|V|=s^2<16N^2$.

Fix distinct $u,v$ and requested colors $\mu,\nu$.
If $\mu=\nu=\chi(u,v)$, all $s-2\ge2$ other points of the line $uv$
are witnesses. Otherwise, choose a direction from
$\operatorname{Dir}_\mu$ at $u$ and a direction from
$\operatorname{Dir}_\nu$ at $v$, neither equal to the direction of $uv$,
and choose the two directions to be distinct. Their two lines intersect
at a unique witness outside $\{u,v\}$. There are at least two choices:
if the colors differ, their direction sets are disjoint and at most one
set loses a direction to $uv$; if they agree but differ from $\chi(u,v)$,
two directions from their common set can be taken in either order.
Distinct direction pairs give distinct intersection points, since a
point outside $\{u,v\}$ determines both directions uniquely. Thus the
coloring has the two-witness property.
\end{proof}
For $N=2^k$, one can choose $s=2N$, giving $4^{k+1}$ vertices.
This construction supplies selected domain sizes, rather than every large
size; induced subgraphs need not retain the two-witness property.
